\section{The Sexes and Race}

\begin{quotex}
In this and in the following section, \textbf{Julius Evola} deals with the traditional view of the sexes. He will justify the notion that the division into sexes has a corporeal, psychic, and spiritual component. His opinion is quite ``politically incorrect''. Oddly, he had to deal with PC in his own time: the German translation omits all references to Germany or Germanic people.
\end{quotex}


\paragraph{Text}


On the bases of the above-mentioned ideas, it would also be necessary to return to the question of hybrids, bringing the sexes into the discussion. Even here we encounter a curious contradiction in race theory, that such questions are almost never posed. Contradicting this principle, claiming the differences of all human types as subjected in equal measure to the same biological laws that race theory considers in the same way, so it seems never to have been considered that normally the inheritance and the power of race can have a different weight according to whether it is about a man or a woman. Anyone who faced this problem had to resolve it directly in reverse, by supposing, again on the basis of simply biological considerations, a greater power of conservation of the race and the type in the woman.

From the point of view of traditional doctrine, it is exactly the opposite that is true in the case of normal humanity. This doctrine, if worthy of the rather minor attention of what today some of the most theatrical and insignificant grant to biological considerations, would furnish some rather useful suggestions for a problem of no small importance, which is the technique for the elevation of relatively inferior breeds through various cycles of heredity. Thus, in the most ancient Indo-European code, the Laws of Manu (Manava Dharmasastra), the transition of a non-Aryan into the caste of the Aryans is possible after seven generations of crosses maintained on the male line. This number seven also reappears in other traditions in analogous circumstances, while, in regard to the cycle of an individual human life, it is the number of years that are necessary for a periodic rejuvenation of all the elements of the body according to scientific research.

In this situation, the Laws of Manu stipulate what must be considered as a cornerstone for the question from the traditional point of view: male inheritance cannot be put on the same level as the female, because, in principle, the former has the ``dominant'' quality, as it is called in Mendelism, and the latter is ``recessive''. Therefore, when the woman is of the higher race, her higher inheritance is overpowered in the mixture, while the superior male inheritance, in the opposite case, is not necessarily contaminated, except in limit-cases or exceptions, and except what we will say about being man. ``Whatever are the qualities of a man with whom a woman is united by legitimate rite, she acquires them like the water of a river united to the ocean'' (Laws of Manu IX, 22).

And again (IX, 33-36): ``If one compares the creative power of the male with that of the female, the male must be declared superior because the ancestor of all beings is distinct from the male characteristics. Whatever type of seed that is thrown in a field prepared in the appropriate season, this seed is developed in a plant with particular qualities, which are those of the male seed.''

\begin{quotex}
By sacred tradition the woman is declared to be the soil, the man is declared to be the seed; the production of all corporeal beings takes place through the union of the soil with the seed. In some cases the seed is more distinguished, and in some the womb of the female; but when both are equal, the offspring is most highly esteemed. On comparing the seed and the receptacle, the seed is declared to be more important; the offspring of all created beings is marked by the characteristics of the seed. Whatever seed is sown in a field, prepared in due season, a plant of that same kind, marked with the peculiar qualities of the seed, springs up in it.
\end{quotex}


Completing the image, all the more, we can concede that when the field is not prepared and the season is not suitable, the male quality will be obstructed in the descendants or will wilt, or certainly it will dry up, but it could never happen that the seed of a palm might come from a juniper plant through a miraculous power of the soil or season -- i.e., by analogy, of the woman and the psychic conditions of a sexual union. As we have deliberately revealed, as long as we have the normal world in view, that is always presupposed by every traditional teaching.

Thus, in order to know what one must think today in this regard, rather than consulting the biologist, it is necessary to clarify the extent to which the modern world can indeed be called normal in regard to the condition of the sexes. The answer, unfortunately, can only be negative. The modern world no longer knows what it means to be a man or a woman in the higher sense. It is moving toward a lack of differentiation of types that is already very visible on the spiritual plane and, from it, here and there it seems to be translated onto the physical and biological plane giving rise to worrisome phenomena. In the West, masculinity and femininity are considered as things simply of the body, instead of qualities, above all, of the inner being of the soul and the spirit. In this respect, in the West, little more than nothing has been known about the polarity, the gap, the different function and dignity of the two sexes for some time. And so, very important problems regarding race are today considered in their exterior and consequential aspects, rather than in their internal and substantial aspects. For example, they are very concerned with the demographic problem and they create every type of institution for hygiene and social assistance and the increase of the race in the strict sense, but they ignore the fundamental point, which is the meaning of the relationship between the sexes and the precise imperative that whoever is born a man, is a man, and woman, a woman, in everything and for everything, in the spirit and in the body, without mixing and without attenuation. Only in this case the traditional teachings have validity and open up almost unlimited possibilities for the initiatives of selection and elevation of the races through adequate crosses and hereditary processes. Certainly not in the case in which, like today, we see in regard to being man or woman, a still more oblique mixing than in regard to being of one race or another: in which some beings are men in the body but feminine in the soul or in the spirit, and vice versa, while being silent about the spread of sexual and psychical inclinations of a really pathological character.

But at this point we have to remind the reader of what we already wrote about it in Revolt against the Modern World, while touching on the death of races. Since descendants are not formed through combinations of hereditary elements made in a laboratory or in appropriate state institutions, but derive from the unions of men and women, it would be logical that, as the premise to every active conception of race and every discrimination of one race or another, the race of the male and the race of the female be defined and separated, in the same corporeal, psychic, and spiritual completeness, by which we formulated the theory of the three levels of race theory.

Moreover, we note a unique circumstance which confirms the fact that races, which have biologically preserved more of the Nordic type, are sometimes found inwardly at a greater level of involution and degeneration than others of the same family: we mean that some of the Nordic peoples---Germans and Anglo-Saxons---are those in which the traditional relationships between the two sexes have been greatly subverted. So-called woman's emancipation---which only means her mutilation and degradation in reality---has in fact begun from those peoples and has had the greatest hold in them, where in the Roman peoples, even if only due to bourgeois and conventional reflexes, something of the normal and traditional way of relating is still preserved, in this regard. The height then is reached when some foreign theorists exalt the triviality of terms like \emph{compagno} and \emph{compagna} and so-called ``respect for woman'' as an alleged characteristic of the Nordic race, not suspecting in the least that by doing so, they simply echo an abnormal state of a relatively recent date, while they assume every conception based on the proper distance, polarity, and different dignity of the two sexes as the Asiatic prejudices of the inferior races of the South. It is necessary to recognize that if such distortions are assumed as principles, the path taken would lead not to the reawakening and the reintegration of the pure Nordic type, but to a further involution---in the direction of a trivialization of an inner leveling of the types---of what still remains in the Germanic peoples.


\flrightit{Posted on 2015-04-23 by Aeneas}

\begin{center}* * *\end{center}

\begin{footnotesize}
\begin{sffamily}
\texttt{Cologero on 2015-04-24 at 08:03 said:}

\textit{Feminism in Norway}\footnote{\url{http://en.wikipedia.org/wiki/Feminism_in_Norway}} confirms what Evola wrote in regard to the feminist movement in the Scandinavian countries.

\hfill

\texttt{argusandphoenix on 2015-04-24 at 21:23 said:}

It's interesting that the British Empire had its two ``greatest'' sovereigns, Queen Elizabeth \& Queen Victoria, as the apex of their monarchy. The males seem to have been largely forgotten or vilified. I remember coming across a reference to this, vis a vis Evola and ``race'', and will try to find it. Contrast this with the Middle Ages, in which ``Good King Richard'' or Henry II are far more important figures. The ancient sources I'm familiar with do indicate that the Teutonic races, when encountered by Rome, had differing conceptions of male-female relationships. I think Caesar remarked on the power and influence held by women among the Celts, as well. It's troubling that the roots of this stretch back almost to pre-history. For what it's worth, Caesar in the Gallic Wars comes across as mild and even-handed and ``frank'' by comparison with the perfidious Celtic tribes which were migrating around Western Europe.

\hfill

\texttt{Mark Citadel on 2015-04-25 at 15:53 said:}

``In this situation, the Laws of Manu stipulate what must be considered as a cornerstone for the question from the traditional point of view: male inheritance cannot be put on the same level as the female, because, in principle, the former has the ``dominant'' quality, as it is called in Mendelism, and the latter is ``recessive''. Therefore, when the woman is of the higher race, her higher inheritance is overpowered in the mixture, while the superior male inheritance, in the opposite case, is not necessarily contaminated, except in limit-cases or exceptions, and except what we will say about being man.''

Fascinating. So even during miscegenation, the male genes overpower regardless of race. This presents an interesting study for the various combinations of miscegenaic birth.

@Cologero -- Say it ain't so that you're actually throwing in the towel on Gornahoor! May I ask why? Have you simply run out of material to work with, or just a general moving on. I've found your analysis of Evola's work to be invaluable to my own.

\hfill

\texttt{Cologero on 2015-04-26 at 00:28 said:}

Mark, I don't believe Evola is referring to biological heredity through ``genes'', but rather the inherited soul and spirit of the race. The next section of Sintesi might expand on that idea.

There is no lack of material, quite the contrary. I'll make clear soon why it is necessary to end one phase and begin another. Moreover, the actual personal work with the Gnosis group is much more important, so it needs to be expanded.

\hfill

\texttt{Nexist418 on 2015-04-26 at 23:55 said:}

Will these be collected together? Were they released in order? If not, will a table of contents be provided?

\hfill

\texttt{Cologero on 2015-04-27 at 07:54 said:}

Yes, Nexist418, the chapters will be collected together and published as a book. They are not being released in order, and there is a table of contents online if you bother to look for it.

Can you tell us what value this work has for you?

\hfill

\texttt{Nexist418 on 2015-04-27 at 13:31 said:}

I wanted to read it in order. I did not find the TOC when I looked, but I will look again.

\hfill

\texttt{Lu Dongbin on 2015-04-28 at 09:36 said:}

Strange situation regarding the Germanic relationship between man and women and the overall masculinity and heroism of the Germanic people in relation to that of Southern Europe, Italians, etc. which Evola himself commented on in many places. From his writing on Disraeli and the British Empire:

``England possesses a monarchy, an almost feudal nobility, and a military caste which, at least up until very recent years, showed remarkable qualities of character and of sang-froid.''

Or from Seyyed Hossein Nasr recounting his meeting with Evola:

``While he extolled the ancient Romans and their virtues, he spoke pejoratively about his contemporary Italians. When I asked him what happened to those Roman virtues, he said they traveled north to Germany and we were left with Italian waiters singing o sole mio!''

So how is it that the masculinity, heroism, and Roman virtues Evola admired are at once concentrated to a greater degree among the Germanic people since the Fall of Rome, and yet these same people have a degenerated relationship between the sexes in comparison? What would be spiritual source or cause of this odd state of affairs?

It seems that Evola was certainly correct about his overall view on the matter, especially regarding ``further involution'' when we take a look at the current state of affairs in the Western world today and especially in Germanic and Nordic countries regarding feminism, etc.

\end{sffamily}
\end{footnotesize}

